\scanpage{58}
\[
\deleted{\unsure{X_1}} \qquad\qquad B\qquad A
\]

\[
X\qquad X\times_S X\supset R
\]

\diagram{\begin{tikzcd}[ampersand replacement=\&,column sep=small]
\& X \arrow[r] \& Y\\
X\times_{\unsure{Y|y}}X\arrow[ur]\arrow[r] \& X\arrow[ur] \&
\end{tikzcd}}

\[
\begin{aligned}
&&x_i\\
S&& B\underset{A}{\otimes}B
\end{aligned}
\]
\[
\unsure{1|\mid}
\]

\diagram{\begin{tikzcd}[ampersand replacement=\&, row sep=small]
B\arrow[d,dash] \& A\arrow[r,dash] \& B\otimes_A B\\
A \& \& B
\end{tikzcd}}

\[
x_i\otimes 1\quad\text{base de}\quad B\otimes B\quad\text{sur}\quad 1\otimes B
\]
\margin{$/$ \unsure{divisant|diviser} par $\unsure{m|\mathfrak m}\subset A$}

\[
\dot x_i\otimes 1\qquad\mathord{\text{---}}\quad L\otimes_{\deleted{\unsure{A}}k}L\quad\text{sur}\quad1\otimes L
\]

$x_i$ base de $L$ sur $k$

$x_i$ \unsure{générateurs} de $B$ sur $A$ par Nakayama.

\[
\sum\lambda_i x_i=0\qquad\lambda_i\in A
\]

\[
\deleted{\sum\lambda_i x_i}\left(\sum\lambda_i x_i\right)\otimes_A1
=\sum(x_i\otimes1)(1\otimes\lambda_i)=0
\]

\[
1\otimes\lambda_i=0\quad\text{pour }\text{tt}\quad i
\]

\unsure{Utilisant|Appliquant} $B\otimes_A B\unsure{\mathbin{-}|\to}B$, on trouve.

\[
\deleted{\lambda_i\otimes1=0\qquad1\otimes\lambda_i=\lambda_i\otimes1}
\]

\[
\lambda_i=0\quad\text{dans}\quad B
\]
% archive sheet number: 57
